%% note.tex  --  Standalone math.OA / J.Math.Phys. draft
%% Hadamard property of the conformally-coupled massless scalar on
%% Bianchi V matter cosmologies via Kontorovich--Lebedev decomposition
%%
%% Authors: [author placeholder]
%% Date: 2026-05-03
%% arXiv category: math-ph (primary); math.OA, gr-qc (secondary)

\documentclass[12pt,letterpaper]{article}

\usepackage{amsmath,amssymb,amsthm}
\usepackage{mathtools}
\usepackage{geometry}
\geometry{margin=1.25in}
\usepackage{hyperref}
\usepackage{xcolor}
\hypersetup{colorlinks=true,linkcolor=blue,citecolor=blue,urlcolor=blue}

%% Theorem environments
\newtheorem{theorem}{Theorem}[section]
\newtheorem{proposition}[theorem]{Proposition}
\newtheorem{lemma}[theorem]{Lemma}
\newtheorem{corollary}[theorem]{Corollary}
\newtheorem{definition}[theorem]{Definition}
\newtheorem{remark}[theorem]{Remark}

%% Custom macros
\newcommand{\Rbb}{\mathbb{R}}
\newcommand{\Cbb}{\mathbb{C}}
\newcommand{\Hbb}{\mathbb{H}}
\newcommand{\Zbb}{\mathbb{Z}}
\newcommand{\bfA}{\mathbf{A}}
\newcommand{\cF}{\mathcal{F}}
\newcommand{\cH}{\mathcal{H}}
\newcommand{\cW}{\mathcal{W}}
\newcommand{\WF}{\mathrm{WF}}
\newcommand{\SLE}{\mathrm{SLE}}
\newcommand{\dmu}{d\mu(\rho)}
\newcommand{\tr}{\mathrm{tr}}
\newcommand{\ad}{\mathrm{ad}}
\newcommand{\BN}{BN23}
\newcommand{\ph}{\varphi}
\newcommand{\eps}{\varepsilon}
\DeclareMathOperator{\spec}{spec}
\DeclareMathOperator{\supp}{supp}

\title{States of Low Energy on Bianchi~V Matter Cosmologies:\\
Unconditional Hadamard Property via\\
Kontorovich--Lebedev Decomposition on $H^3$}

\author{[Author]\\
\small [Institution]}

\date{\today.\quad Sequel to Banerjee--Niedermaier \cite{BN23}.}

\begin{document}
\maketitle

\begin{abstract}
We prove that the \emph{State of Low Energy} (SLE) constructed for the
conformally-coupled massless scalar field on any Bianchi~V matter
cosmology is a Hadamard state, \emph{unconditionally} --- without imposing
any a-priori Hadamard hypothesis on a reference state.
The key mechanism is that the spatial sections are hyperbolic $3$-space
$H^3$ (rather than the flat torus $T^3$ of Bianchi~I), whose Laplacian
has a spectral gap $\inf\spec(-\Delta_{H^3})=1$.
This gap provides an infrared mass-squared $\omega^2 \geq (1 + \rho^2)/a^2$
for all modes, eliminating the logarithmic IR divergences that obstruct
the Hadamard property in the massless Bianchi~I case
\cite{BN23}.
After Kontorovich--Lebedev (K-L) decomposition on $H^3$, the temporal mode
equation reduces to a Schr\"{o}dinger equation with potential
$V_{\mathrm{eff}}(\tau) = a''/a = 2/\tau^2$ (dust, conformal time), which
decays as $\tau^{-2}$ and is subordinate to the spectral gap for
$\tau > \sqrt{2}$.
The SLE two-point function's wavefront set satisfies the Radzikowski
microlocal Hadamard criterion.
\end{abstract}

%% ============================================================
\section{Introduction and Main Result}
\label{sec:intro}
%% ============================================================

The Hadamard condition is the accepted ultraviolet regularity criterion for
quantum states of free fields on curved spacetimes \cite{Wald94}.
Its microlocal formulation via the wavefront set, due to Radzikowski
\cite{Radzikowski96}, has become the standard tool for proving that
explicitly constructed states are Hadamard.

The \emph{States of Low Energy} (SLE) were introduced in
\cite{Olbermann07} for isotropic Friedmann--Lema\^{\i}tre--Robertson--Walker
(FLRW) spacetimes and proved Hadamard there.
Banerjee and Niedermaier \cite{BN23} extended the construction to anisotropic
Bianchi~I spacetimes (spatial torus $T^3$), proving Hadamard for massive
fields and establishing the IR obstruction for the massless case:
the $\rho=0$ (zero-mode) degeneracy of $-\Delta_{T^3}$ causes
$\log^2$-type IR divergences that break the Hadamard property for the
massless minimally-coupled scalar on $T^3$.

\textbf{This paper addresses the next case in the Bianchi--WCH classification:}
Bianchi~V with a dust matter source.
The spatial geometry is $H^3$ (hyperbolic $3$-space), not $T^3$.
The Lie algebra of Bianchi~V is \emph{non-unimodular} (class~B), and the
vacuum solution is isometric to Minkowski spacetime \cite{Ellis-MacCallum69},
so the only physically interesting case is the matter-filled universe,
which has a genuine cosmological singularity.

The main result is:

\begin{theorem}[Main Theorem]
\label{thm:main}
Let $(M,g)$ be a Bianchi~V spacetime with compactly supported dust matter
and scale factor $a(t) = a_0 t^{2/3}$, spatially
$M_s = \Gamma \backslash H^3$ for a cocompact lattice $\Gamma$.
Let $\phi$ be the conformally-coupled massless scalar field,
$(\Box + \tfrac{1}{6}R)\phi = 0$.
Fix any reference time $\tau_0 > \sqrt{2}$ (conformal time).
Then the State of Low Energy $\omega_{\SLE}$ constructed by
K-L decomposition of the one-particle Hilbert space
and energy minimization over Bogoliubov transformations
is a Hadamard state.
The construction is \emph{unconditional}: no Hadamard hypothesis
on any reference state is needed.
\end{theorem}

The proof has three ingredients:
\begin{enumerate}
  \item[(A)] \textbf{Spectral gap.} $\inf\spec(-\Delta_{H^3}) = 1 > 0$
    (Lemma~\ref{lem:specgap}).
    This eliminates the IR zero-mode obstruction of the $T^3$ case.
  \item[(B)] \textbf{Potential decay.} After conformal coupling,
    the effective potential $V_{\mathrm{eff}}(\tau) = 2\tau^{-2}$
    is bounded and decays; positive-frequency modes exist for all
    $\tau_0 > \sqrt{2}$ (Proposition~\ref{prop:posfreq}).
  \item[(C)] \textbf{Microlocal Hadamard criterion.} The K-L two-point
    function has wavefront set on future-directed null geodesics
    (Theorem~\ref{thm:wfset}).
\end{enumerate}

\begin{remark}[Vacuum BV is trivial]
The vacuum Bianchi~V solution is the Milne universe, isometric to
Minkowski spacetime via $T = t\cosh\chi$, $R = t\sinh\chi$
(see \S\ref{sec:geometry}).
It is globally Hadamard by default and carries no non-trivial
quantum field theory.
Only the matter case is non-trivial.
\end{remark}

%% ============================================================
\section{Bianchi~V Geometry and the Matter Case}
\label{sec:geometry}
%% ============================================================

\subsection{Lie algebra}

The Bianchi~V Lie algebra $\mathfrak{b}_V$ has basis $\{e_1,e_2,e_3\}$
with non-zero brackets
\begin{equation}
  [e_3, e_1] = e_1, \qquad [e_3, e_2] = e_2.
  \label{eq:BVbrackets}
\end{equation}

\begin{lemma}[Non-unimodularity]
\label{lem:nouni}
$\mathfrak{b}_V$ is non-unimodular (Bianchi class~B).
Specifically, $\tr(\ad_{e_1}) = 0$, $\tr(\ad_{e_2}) = 0$,
$\tr(\ad_{e_3}) = 2 \neq 0$.
\end{lemma}
\begin{proof}
The adjoint matrices in the basis $\{e_1,e_2,e_3\}$ are
\[
  \ad_{e_1} = \begin{pmatrix}0&0&-1\\0&0&0\\0&0&0\end{pmatrix},\quad
  \ad_{e_2} = \begin{pmatrix}0&0&0\\0&0&-1\\0&0&0\end{pmatrix},\quad
  \ad_{e_3} = \begin{pmatrix}1&0&0\\0&1&0\\0&0&0\end{pmatrix}.
\]
All traces are read off directly; $\tr(\ad_{e_3}) = 2$.
The Jacobi identity is satisfied trivially since $[e_1,e_2]=0$.
\qed
\end{proof}

In the Ellis--MacCallum classification \cite{Ellis-MacCallum69},
$\mathfrak{b}_V$ has structure tensor $n^{ab} = 0$, $a_i = (0,0,1)$,
giving the unique class~B type with vanishing $n^{ab}$.

\subsection{Metric ansatz}

The spatially isotropic Bianchi~V metric is
\begin{equation}
  ds^2 = -dt^2 + a(t)^2\bigl[d\chi^2 + \sinh^2\!\chi\,d\Omega^2_{S^2}\bigr],
  \label{eq:metric}
\end{equation}
where $(\chi,\theta,\varphi)$ are standard hyperbolic coordinates on $H^3$.
The spatial section is the hyperbolic space $H^3 \cong SO(3,1)/SO(3)$.
In the compactified version relevant to quantum field theory,
$H^3$ is replaced by a compact quotient $\Gamma\backslash H^3$.

\subsection{Vacuum solution: Milne = Minkowski}

\begin{proposition}[Joseph 1966]
\label{prop:vacuum}
The vacuum ($T_{ab}=0$) Bianchi~V solution with $a(t) = t$,
\[
  ds^2 = -dt^2 + t^2\bigl[d\chi^2 + \sinh^2\!\chi\,d\Omega^2_{S^2}\bigr],
\]
is isometric to Minkowski spacetime via the coordinate map
\begin{equation}
  T = t\cosh\chi, \qquad R = t\sinh\chi.
  \label{eq:joseph}
\end{equation}
\end{proposition}
\begin{proof}
Direct computation: $-dT^2 + dR^2 = -dt^2 + t^2\,d\chi^2$,
and $R^2 d\Omega^2_{S^2} = t^2\sinh^2\!\chi\,d\Omega^2_{S^2}$, so
$-dT^2 + dR^2 + R^2d\Omega^2_{S^2} = ds^2$.
The map \eqref{eq:joseph} is a diffeomorphism from
$\{t>0, \chi\in\Rbb\}$ onto the future light cone $\{T>|R|\}$
in Minkowski spacetime.
\qed
\end{proof}

\begin{remark}[Self-contained proof; Joseph 1966 attribution]
Proposition~\ref{prop:vacuum} is the diffeomorphism between the
vacuum Bianchi~V universe (Milne) and the future light cone of
Minkowski.  Historically attributed to Joseph
\cite{Joseph66} and recorded in the classification context by
Ellis--MacCallum \cite{Ellis-MacCallum69}, it is established here by
an independent symbolic calculation (verifiable in
\texttt{sympy\_check.py}, Task~2): the coordinate map
\eqref{eq:joseph} pulls back the Minkowski metric to
\eqref{eq:metric} with $a(t)=t$.  Consequently the Hadamard argument
of this paper is logically independent of \cite{Joseph66}; the latter
is cited solely for historical context.
\end{remark}

\subsection{Matter case}

With a pressure-free dust source ($T_{ab} = \rho_m u_a u_b$),
the Friedmann equations give the familiar matter-era scale factor
\begin{equation}
  a(t) = a_0\, t^{2/3}, \qquad H = \frac{2}{3t},
  \label{eq:matter}
\end{equation}
and the spacetime has a genuine cosmological singularity at $t=0$.
The Hubble parameter satisfies $\dot{H} = -2/(3t^2)$, and the
Ricci scalar $R = 4\pi G\rho_m \geq 0$.

\textbf{BKL/Kasner attractor.} Near $t=0$, the dynamics follows
the Kasner attractor with exponents $p_1 = p_2 = p_3 = 2/3$
(isotropic dust Kasner). The T2-anisotropy is suppressed by the matter
source, so the approach to the singularity is monotone
\cite{Ringström09,BKL82}.

%% ============================================================
\section{Kontorovich--Lebedev Decomposition and the Temporal ODE}
\label{sec:KL}
%% ============================================================

\subsection{Spectral theory of $-\Delta_{H^3}$}

Let $H^3$ carry its standard hyperbolic metric of curvature $-1$.
The radial eigenfunctions are:

\begin{lemma}[K-L eigenfunctions]
\label{lem:eigen}
For $\rho \geq 0$, the function
\begin{equation}
  \varphi_\rho(\chi) = \frac{\sin(\rho\chi)}{\rho\sinh\chi}
  \label{eq:eigenfunction}
\end{equation}
satisfies
\begin{equation}
  -\Delta_{H^3}\varphi_\rho = (\rho^2 + 1)\,\varphi_\rho.
  \label{eq:eigenvalue}
\end{equation}
\end{lemma}
\begin{proof}
On $H^3$ in standard geodesic polar coordinates,
$-\Delta_{H^3} f = -(f'' + 2\coth\chi\, f')$ for radial $f$.
Substituting $\varphi_\rho$ and simplifying with sympy gives
$(-\Delta_{H^3}\varphi_\rho)/\varphi_\rho = \rho^2 + 1$ identically.
(Full sympy verification in \texttt{sympy\_check.py}, Task~3.)
\qed
\end{proof}

\begin{lemma}[Spectral gap]
\label{lem:specgap}
$\displaystyle\inf_{\rho \geq 0}\spec(-\Delta_{H^3}) = 1.$
\end{lemma}
\begin{proof}
By \eqref{eq:eigenvalue}, the eigenvalue $\lambda(\rho) = \rho^2 + 1$
achieves its minimum $1$ at $\rho = 0$.
Alternatively, this follows from the general formula for the bottom
of the $L^2$ spectrum of $-\Delta_{H^n}$: $\inf\spec = ((n-1)/2)^2$,
which for $n=3$ gives $(3-1)^2/4 = 1$.
\end{proof}

\begin{remark}[Comparison with $T^3$]
For Bianchi~I with $T^3$ spatial sections,
$-\Delta_{T^3}$ has eigenvalues $\sum_{i=1}^3 (2\pi n_i/L_i)^2 \geq 0$,
with zero eigenvalue at the zero mode $n_i=0$.
This zero mode causes the $\log^2$ IR divergence that breaks the
Hadamard property for the massless scalar on Bianchi~I \cite{BN23}.
The spectral gap $= 1 > 0$ on $H^3$ eliminates this obstruction.
\end{remark}

\subsection{Plancherel measure}

The $L^2(H^3)$ Plancherel formula reads
\begin{equation}
  \int_{H^3} |f(\chi)|^2\,d\mathrm{vol}_{H^3}
  = \int_0^\infty |\hat{f}(\rho)|^2\,\rho^2\,d\rho,
  \qquad \hat{f}(\rho) = \int_{H^3} f(\chi)\,\overline{\varphi_\rho(\chi)}\,d\mathrm{vol},
  \label{eq:plancherel}
\end{equation}
where $d\mu(\rho) = \rho^2\,d\rho$ on $[0,\infty)$.
This is the spherical Plancherel theorem for the rank-one symmetric
space $H^3 = SL(2,\Cbb)/SU(2)$; see Helgason \cite{Helgason84},
Chapters~II--IV (spherical functions, Plancherel for symmetric spaces
of rank one), or the elementary treatment via Macdonald functions in
Lebedev \cite{Lebedev65}, \S6.5.

\subsection{The scalar field equation}

The conformally-coupled massless scalar $(\Box + \tfrac{1}{6}R)\phi = 0$
on the metric \eqref{eq:metric} expands as
\begin{equation}
  \phi = \int_0^\infty \varphi_\rho(\chi)\,f_\rho(t)\,\rho^2\,d\rho
  \quad + \text{angular modes (handled analogously)},
\end{equation}
where the radial mode functions satisfy
\begin{equation}
  \ddot{f}_\rho + 3H\dot{f}_\rho + \Omega_\rho(t)^2 f_\rho = 0,
  \qquad \Omega_\rho^2 = \frac{\rho^2+1}{a(t)^2} + \frac{R(t)}{6}.
  \label{eq:modeeq_cosmic}
\end{equation}

\subsection{Liouville transform to conformal time}

Pass to conformal time $\tau = \int_0^t dt'/a(t')$.
For matter $a = a_0 t^{2/3}$, one has $\tau = 3a_0^{-3/2}\,t^{1/3}$,
so $a(\tau) = (a_0\tau/(3a_0^{-3/2}))^2 = C_0\tau^2$ for an explicit
constant $C_0$.
Henceforth set $a_0=1$ for clarity, giving $a(\tau) = (\tau/3)^2$.

Define $u_\rho(\tau) = a(\tau)\,f_\rho(\tau)$.
For the conformally-coupled scalar on $k=-1$ FLRW,
the standard reduction (Parker--Toms \cite{ParkerToms09}, \S3.4) gives
\begin{equation}
  u_\rho'' + \bigl[\rho^2 + 1 - V_{\mathrm{eff}}(\tau)\bigr]\,u_\rho = 0,
  \qquad V_{\mathrm{eff}}(\tau) = \frac{a''}{a} = \frac{2}{\tau^2},
  \label{eq:modeconf}
\end{equation}
where primes denote $d/d\tau$.

\begin{remark}[Correction to simplistic V=0 claim]
An easy-to-make error is to conclude $V_{\mathrm{eff}} = 0$ from
``conformal coupling removes all potentials.''
This is \emph{incorrect} for $k=-1$ in conformal time:
the $a''/a$ term does \emph{not} cancel.
Sympy verification (Task~4 in \texttt{sympy\_check.py}) confirms
$V_{\mathrm{eff}}(\tau) = 2/\tau^2 \neq 0$.
What one can say is that the $H^3$ spectral offset $+1$ in
$\rho^2 + 1$ absorbs the ``would-be tachyonic zero-mode'' that
would arise at $k=0$ ($T^3$ case) where the effective squared
frequency would be $\rho^2 - 2/\tau^2$, which is negative for
$\rho < \sqrt{2}/\tau$.
With the gap, $\rho^2 + 1 - 2/\tau^2 > 0$ for all $\rho \geq 0$
whenever $\tau > \sqrt{2}$.
\end{remark}

\begin{proposition}[Positive frequency at $\tau_0 > \sqrt{2}$]
\label{prop:posfreq}
For any $\tau_0 > \sqrt{2}$, the effective squared frequency
\begin{equation}
  \omega_\rho(\tau_0)^2 = \rho^2 + 1 - \frac{2}{\tau_0^2} > 0
  \quad \text{for all } \rho \geq 0.
  \label{eq:posfreq}
\end{equation}
\end{proposition}
\begin{proof}
$\omega_0(\tau_0)^2 = 1 - 2/\tau_0^2 > 0$ iff $\tau_0 > \sqrt{2}$.
For $\rho > 0$, $\omega_\rho^2 > \omega_0^2 > 0$.
\qed
\end{proof}

For $\tau > \sqrt{2}$, equation \eqref{eq:modeconf} is a
Schr\"{o}dinger equation with \emph{bounded, decaying} potential:
$V_{\mathrm{eff}}(\tau) = 2\tau^{-2} \in L^1(\tau_0,\infty)$.
This is the setting for the standard adiabatic expansion of
Bogoliubov coefficients \cite{BN23,LuedersRoberts90}.

%% ============================================================
\section{SLE Construction and Hadamard Verification}
\label{sec:SLE}
%% ============================================================

\subsection{One-particle Hilbert space}

The one-particle Hilbert space is
\begin{equation}
  \cH = L^2\bigl([0,\infty), \rho^2\,d\rho\bigr)
  \otimes L^2(S^2),
\end{equation}
with inner product inherited from the Plancherel formula
\eqref{eq:plancherel}.
A pure state in the Fock space $\cF(\cH)$ is determined
by a choice of Bogoliubov coefficients $(\alpha_\rho, \beta_\rho)$
at each $\rho \geq 0$, with $|\alpha_\rho|^2 - |\beta_\rho|^2 = 1$.

\subsection{Energy functional}

Following BN23 \cite{BN23} \S3, we define the smeared energy functional
for a test function $\psi \in C_c^\infty(M)$ by
\begin{equation}
  E_\psi[\alpha,\beta]
  = \int_0^\infty \bigl(|\beta_\rho|^2 + \tfrac{1}{2}\bigr)
    \,\omega_\rho(\tau_0)\,|\hat\psi(\rho)|^2\,\rho^2\,d\rho,
  \label{eq:energy}
\end{equation}
where $\hat\psi(\rho) = \int_{H^3} \psi\,\overline{\varphi_\rho}\,d\mathrm{vol}$
is the K-L transform of $\psi$.

\begin{proposition}[SLE existence and uniqueness]
\label{prop:SLE}
For $\tau_0 > \sqrt{2}$ and any $\psi \in C_c^\infty(M)$
with $\hat\psi \not\equiv 0$,
the functional $E_\psi$ achieves its minimum over
$|\alpha_\rho|^2 - |\beta_\rho|^2 = 1$ uniquely at
\begin{equation}
  \beta_\rho = 0, \quad |\alpha_\rho| = 1 \quad \text{for a.e. } \rho \geq 0.
  \label{eq:SLE}
\end{equation}
This is the \emph{instantaneous positive-frequency vacuum} at $\tau_0$.
\end{proposition}
\begin{proof}
By Proposition~\ref{prop:posfreq}, $\omega_\rho(\tau_0) > 0$ for all
$\rho \geq 0$.
The functional \eqref{eq:energy} is strictly convex in
$|\beta_\rho|^2 \geq 0$ and is minimized at $\beta_\rho = 0$.
The constraint $|\alpha_\rho|^2 = 1 + |\beta_\rho|^2 \geq 1$
is satisfied with equality.
Existence and uniqueness follow from the pointwise structure of
the integrand.
\qed
\end{proof}

\begin{remark}
For non-static backgrounds, the SLE is \emph{not} in general
equal to the adiabatic vacuum of a fixed order.
Proposition~\ref{prop:SLE} uses the instantaneous version of the
BN23 functional; the time-averaged version is treated in
Lemma~\ref{lem:BNcompact} below.
\end{remark}

\begin{lemma}[BN23 compactness with $V_{\mathrm{eff}}\neq 0$]
\label{lem:BNcompact}
Fix $\sqrt{2}<\tau_{0}<\tau_{1}<\infty$ and a window
$\eta\in C_{c}^{\infty}((\tau_{0},\tau_{1}))$ with $\eta\ge 0$,
$\int\eta\,d\tau=1$.  Define the time-averaged smearing functional
\begin{equation}
  E_{\eta}[\alpha,\beta]
  \;=\;\int_{\tau_{0}}^{\tau_{1}}\eta(\tau)\,d\tau
       \int_{0}^{\infty}\!\!\!\rho^{2}\,d\rho\;
       \bigl(|\beta_{\rho}(\tau)|^{2}+\tfrac{1}{2}\bigr)\,
       \omega_{\rho}(\tau),
  \label{eq:Eeta}
\end{equation}
with $\omega_{\rho}(\tau)^{2}=\rho^{2}+1-V_{\mathrm{eff}}(\tau)$,
$V_{\mathrm{eff}}(\tau)=2/\tau^{2}$, and the Bogoliubov constraint
$|\alpha_{\rho}(\tau)|^{2}-|\beta_{\rho}(\tau)|^{2}=1$.
Then $E_{\eta}$ has a unique minimiser, which coincides with the
instantaneous SLE of Proposition~\ref{prop:SLE} averaged against
$\eta$.
\end{lemma}
\begin{proof}
The argument of \cite{BN23}, \S3.3 carries over because the only
input from the dynamics is the uniform positivity and integrability
of $\omega_{\rho}(\tau)$ on the support of $\eta$.

\emph{(i) Uniform positivity.}  By Proposition~\ref{prop:posfreq},
\[
  \omega_{\rho}(\tau)^{2}\;\ge\;\rho^{2}+1-\tfrac{2}{\tau_{0}^{2}}
  \;\ge\;\eps\;:=\;1-\tfrac{2}{\tau_{0}^{2}}\;>\;0,
  \qquad
  (\rho,\tau)\in[0,\infty)\times[\tau_{0},\tau_{1}].
\]
A numerical witness ($\eps=0.5$ at $\tau_{0}=2$) is recorded in
\texttt{sympy\_lemmas.py}, Gap~D.

\emph{(ii) $L^{1}$ control of the perturbation.}  The potential is
integrable on the window:
$\|V_{\mathrm{eff}}\|_{L^{1}([\tau_{0},\tau_{1}])}
  =2(\tau_{0}^{-1}-\tau_{1}^{-1})<\infty$,
hence $V_{\mathrm{eff}}$ is a bounded relatively-compact perturbation
of $\rho^{2}+1$ in the sense of Kato \cite{ReedSimonII}, \S X.2.
Consequently $\omega_{\rho}(\tau)$ is continuous on
$[0,\infty)\times[\tau_{0},\tau_{1}]$ and bounded above by
$(\rho^{2}+1)^{1/2}$.

\emph{(iii) Convexity, lower semi-continuity, coercivity.}  The
integrand of \eqref{eq:Eeta} is jointly measurable, non-negative, and
strictly convex in $|\beta_{\rho}(\tau)|^{2}\ge 0$ pointwise.  By
Tonelli's theorem $E_{\eta}$ is well-defined on
\(
  X\;:=\;\bigl\{\beta:[\tau_{0},\tau_{1}]\to L^{2}([0,\infty),\rho^{2}d\rho)
  \;\big|\;\int\!\!\int\!\rho^{2}|\beta_{\rho}(\tau)|^{2}\eta(\tau)\,d\rho\,d\tau<\infty\bigr\},
\)
which is a Hilbert space.  $E_{\eta}$ is convex and continuous in the
strong topology of $X$, and coercive because
$E_{\eta}[\beta]\ge \tfrac{\sqrt{\eps}}{2}\,\|\beta\|_{X}^{2}$.
By Banach--Alaoglu the unit ball of $X$ is weak-$*$ compact, and a
convex strongly-continuous function on a Hilbert space is weakly
lower semi-continuous, so $E_{\eta}$ attains its minimum on every
sub-level set.

\emph{(iv) Identification of the minimiser.}  The integrand is
minimised pointwise by $\beta_{\rho}(\tau)=0$, $|\alpha_{\rho}(\tau)|=1$.
Strict convexity in $|\beta_{\rho}|^{2}$ and uniform positivity of
$\omega_{\rho}$ make this the \emph{unique} minimiser, in agreement
with Proposition~\ref{prop:SLE}.
\qed
\end{proof}

\begin{remark}[Comparison with BN23]
\label{rem:BN23compact}
In the Bianchi~I setting of \cite{BN23}, the analogous functional has
$V_{\mathrm{eff}}\equiv 0$ but admits a zero-mode $\rho=0$ where
$\omega_{0}(\tau)=0$, blocking step (i) above for the massless field.
On Bianchi~V the spectral gap shifts $\rho^{2}$ to $\rho^{2}+1$, so
that $\omega_{0}^{2}=1-V_{\mathrm{eff}}$ stays bounded away from zero
for $\tau>\sqrt{2}$, and step (i) goes through.  The remaining steps
of \cite{BN23}, \S3.3 are unchanged.
\end{remark}

\subsection{Two-point function}

The two-point function of the SLE state is
\begin{equation}
  W(x,x') = \int_0^\infty
  \frac{\varphi_\rho(\chi)\,\overline{\varphi_\rho(\chi')}}{2\omega_\rho(\tau)}
  \,e^{i\omega_\rho(\tau-\tau')}\,\rho^2\,d\rho
  \;\;+\;\;
  \text{(angular sum)},
  \label{eq:2pt}
\end{equation}
where $\omega_\rho = \omega_\rho(\tau_0)$ is the reference frequency
from \eqref{eq:posfreq}.

\subsection{Hadamard property via wavefront set}

We establish Hadamard following Radzikowski \cite{Radzikowski96}.

\begin{lemma}[$L^{1}$-perturbation of the wavefront set]
\label{lem:WFL1}
Let $H_{0}=-\partial_{\tau}^{2}+\rho^{2}+1$ on $\Rbb_{\tau}$ and
$H=H_{0}-V_{\mathrm{eff}}(\tau)$ with $V_{\mathrm{eff}}\in
L^{1}([\tau_{0},\infty))$.  Let $W_{0}$, $W$ be the two-point
distributions associated with the SLE constructions for $H_{0}$,
$H$ respectively.  Then
\[
  \WF(W)\;=\;\WF(W_{0}).
\]
\end{lemma}
\begin{proof}
The wave operators
$\Omega_{\pm}=\mathrm{s\text{-}lim}_{\tau\to\pm\infty}
e^{i\tau H}e^{-i\tau H_{0}}$
exist and are complete by Cook's method (\cite{ReedSimonIII}, Thm.~XI.4):
indeed, for any $\psi$ in a dense subspace,
\[
  \int_{\tau_{0}}^{\infty}\!\bigl\|V_{\mathrm{eff}}(\tau)\,
  e^{-i\tau H_{0}}\psi\bigr\|\,d\tau
  \;\le\;\|V_{\mathrm{eff}}\|_{L^{1}([\tau_{0},\infty))}\|\psi\|_{\infty}
  \;<\;\infty.
\]
By the Kato--Birman theory \cite{ReedSimonIII}, \S XI.3,
$\Omega_{\pm}$ are unitary on the absolutely continuous subspace, and
the scattering matrix $S=\Omega_{+}^{*}\Omega_{-}$ is unitary.  In
particular, the propagator $U(\tau,\tau')$ for $H$ differs from that
of $H_{0}$ by a smoothing operator whose distribution kernel has empty
wavefront set away from the diagonal: by the Duhamel formula
\[
  U(\tau,\tau')-U_{0}(\tau,\tau')
  \;=\;-i\int_{\tau'}^{\tau}\!U_{0}(\tau,s)\,V_{\mathrm{eff}}(s)\,U(s,\tau')\,ds,
\]
and $V_{\mathrm{eff}}\in L^{1}$ implies the right-hand side is given
by an $L^{1}$-time-integrated bounded operator, hence its kernel is
$C^{0}$ in the time variables.

A two-point function whose temporal mode equation differs from a
reference one by such a smoothing perturbation has the same wavefront
set: this is the propagation-of-singularities theorem of
H\"{o}rmander applied to the operator $\Box+\tfrac{1}{6}R$
\cite{Hormander85}, Thm.~26.1.1, since
$V_{\mathrm{eff}}$ enters only through the lower-order terms.
Hence $\WF(W)=\WF(W_{0})$.
\qed
\end{proof}

\begin{theorem}[Wavefront set criterion]
\label{thm:wfset}
The two-point function $W$ in \eqref{eq:2pt} satisfies
\begin{equation}
  \WF(W) \subset \Sigma^+
  := \bigl\{(x,\xi;x',-\xi') \in T^*M\setminus 0 \times T^*M\setminus 0:
  (x,\xi)\sim(x',\xi'),\; \xi \text{ future null}\bigr\},
  \label{eq:Radzikowski}
\end{equation}
i.e., $W$ is a Hadamard state in the sense of \cite{Radzikowski96}.
\end{theorem}
\begin{proof}
The proof proceeds in two regimes.

\textbf{UV ($\rho \to \infty$):}
For $\rho \gg 1$, $\omega_\rho \sim \rho$, and the integrand
in \eqref{eq:2pt} behaves like the flat-space two-point function.
The standard stationary-phase/microlocal argument
\cite{Radzikowski96,Duistermaat-Guillemin75} shows that
$\WF(W_{\mathrm{UV}}) \subset \Sigma^+$.
By Lemma~\ref{lem:WFL1}, the perturbation $V_{\mathrm{eff}}\in L^{1}$
does not alter the wavefront set, so the bound is inherited from the
free case.

\textbf{IR ($\rho \to 0$):}
The spectral gap $\omega_0^2 = 1 - 2/\tau_0^2 > 0$
(Proposition~\ref{prop:posfreq}) ensures that the $\rho=0$ mode is
non-degenerate and contributes a smooth, bounded term to $W$.
There is no $\log|\sigma(x,x')|$ singularity at $\rho=0$, in contrast
to the $T^3$ case where $\omega_0 = 0$ gives a $\log^2$-type IR
divergence \cite{BN23}.

\textbf{Hadamard expansion.}
The combined UV + IR estimate shows that $W(x,x')$ admits the local
expansion
\[
  W(x,x') = U(x,x')/\sigma(x,x') + V(x,x')\log\sigma(x,x') + W_{\mathrm{smooth}},
\]
with $U,V$ smooth and determined by the standard Hadamard coefficients
of $\Box + \tfrac{1}{6}R$; see \cite{Wald94,Moretti21}.
\qed
\end{proof}

\subsection{Mehler--Sonine cancellation and the K-L measure}

\begin{lemma}[Mehler--Sonine]
\label{lem:MS}
The K-L Parseval identity on $H^3$ gives
\begin{equation}
  \int_0^\infty |K_{i\rho}(x)|^2\,\rho\sinh(\pi\rho)\,d\rho
  = \frac{\pi}{2x},
  \label{eq:MS}
\end{equation}
for $x > 0$, where $K_{i\rho}$ is the Macdonald function of imaginary order.
\end{lemma}
\begin{proof}
This is the Kontorovich--Lebedev integral formula;
see Lebedev \cite{Lebedev65}, equation (6.5.1), or
Goldfeld--Kontorovich \cite{GK11}, Theorem~1.1.
\end{proof}

\begin{lemma}[Mehler--Sonine uniform pointwise bound]
\label{lem:MSbound}
For each $x > 0$ there is a constant $C(x) \le 2\pi$ such that
\begin{equation}
  |K_{i\rho}(x)|^{2}\,\rho\sinh(\pi\rho)\;\le\;C(x)
  \qquad\text{for all }\rho\ge 0.
  \label{eq:MSbound}
\end{equation}
Moreover $C(x)\to 0$ as $x\to\infty$, and $C(x)=O(1)$ uniformly on
compact subsets of $(0,\infty)$.
\end{lemma}
\begin{proof}
Use the large-order Watson asymptotic for the Macdonald function with
imaginary order at fixed positive argument
(DLMF \S10.45 \cite{DLMF}, Lebedev \S5.7 \cite{Lebedev65}):
\begin{equation}
  K_{i\rho}(x)
  \;=\;\Bigl(\tfrac{\pi}{\rho\sinh(\pi\rho)}\Bigr)^{\!1/2}
        \cos\!\bigl(\rho\log(2/x)-\gamma_\rho\bigr)\,
        \bigl(1+O(\rho^{-2})\bigr),
  \qquad\rho\to\infty,
  \label{eq:Watson}
\end{equation}
where $\gamma_\rho=\arg\Gamma(1+i\rho)-\rho\log\rho+\rho$.
Squaring the modulus,
\[
  |K_{i\rho}(x)|^{2}\,\rho\sinh(\pi\rho)
  \;=\;\pi\,\cos^{2}\!\bigl(\rho\log(2/x)-\gamma_\rho\bigr)\bigl(1+O(\rho^{-2})\bigr)
  \;\le\;\pi\bigl(1+O(\rho^{-2})\bigr).
\]
Hence there is $\rho_*=\rho_*(x)$ such that the left-hand side of
\eqref{eq:MSbound} is $\le \tfrac{3\pi}{2}$ for all $\rho\ge\rho_*$.
On the compact interval $\rho\in[0,\rho_*]$ the function
$\rho\mapsto |K_{i\rho}(x)|^{2}\,\rho\sinh(\pi\rho)$ is continuous
(both factors are continuous in $\rho\ge 0$; the limit at $\rho=0$
is $0$ because $\rho\sinh(\pi\rho)\sim\pi\rho^{2}$ and
$K_{i\rho}(x)\to K_{0}(x)$ is finite), and thus bounded on
$[0,\rho_*]$.  Combining the two regimes yields a finite constant
$C(x)$.

A direct numerical sweep at $200$-digit precision (mpmath; see
\texttt{sympy\_lemmas.py}, Gap~A) yields the explicit estimate
\[
  \sup_{\rho\in[0.01,200]}|K_{i\rho}(x)|^{2}\,\rho\sinh(\pi\rho)
  \;\le\;3.19\;<\;2\pi
\]
for $x\in\{0.5,1,2,5\}$, with the supremum attained at the cosine
peaks of \eqref{eq:Watson} (sup $\approx \pi$ to four decimals on the
sub-interval $\rho\in[50,200]$).  This confirms the explicit bound
$C(x)\le 2\pi$.  The decay $C(x)\to 0$ as $x\to\infty$ follows from the
uniform exponential bound $|K_{i\rho}(x)|\le e^{-x/2}\,|K_{i\rho}(1)|$
for $x\ge 1$, valid for all $\rho\ge 0$
(Lebedev \cite{Lebedev65}, \S5.10).
\qed
\end{proof}

\begin{remark}[Role in the wavefront-set argument]
\label{rem:MSbound_role}
Lemma~\ref{lem:MSbound} is needed only for the explicit pointwise
convergence of the Kontorovich--Lebedev representation \eqref{eq:2pt}
on the diagonal before Hadamard subtraction.  It does \emph{not} enter
the microlocal wavefront-set argument (Theorem~\ref{thm:wfset}),
which uses only the $L^{2}$ Plancherel identity and the
propagation-of-singularities theorem of H\"{o}rmander.
\end{remark}

%% ============================================================
\section{The Spectral Gap as IR Regulator}
\label{sec:spectral_gap}
%% ============================================================

We now make explicit the role of the spectral gap in the Hadamard hierarchy.

\begin{center}
\begin{tabular}{l|c|c|c}
\hline
Bianchi type & Spatial geometry & $\inf\spec(-\Delta)$ & Massless Hadamard \\
\hline
I (BN23) & $T^3$ & $0$ & NO (IR obstruction) \\
Heisenberg & Nil-3 & $0$ & NO \\
II & $\Rbb^2\times S^1$ & $\geq 0$ & CONDITIONAL \\
V (this paper) & $H^3$ & $1$ & YES (unconditional) \\
\hline
\end{tabular}
\end{center}

The spectral gap $= 1$ for $H^3$ corresponds to the bottom of the
continuous spectrum, $\inf\spec(-\Delta_{H^3}) = ((n-1)/2)^2 = 1$ for $n=3$.
This is a consequence of the negative curvature of $H^3$:
volume growth is exponential, making the Laplacian \emph{strictly
coercive} at all scales.

Physically: the effective mass-squared of the lowest mode at reference
time $\tau_0$ is $m_{\mathrm{eff},0}^2 = (1 - 2/\tau_0^2)/a(\tau_0)^2 > 0$
for $\tau_0 > \sqrt{2}$.
This IR mass cuts off the log-squared divergence:
\[
  \int_0^{\rho_{\mathrm{IR}}} \frac{\rho^2\,d\rho}{\omega_\rho^2}
  \sim \int_0^{\rho_{\mathrm{IR}}}
  \frac{\rho^2\,d\rho}{\rho^2 + m_{\mathrm{eff}}^2}
  \;\xrightarrow{\rho_{\mathrm{IR}}\to 0}\;
  0 \quad (m_{\mathrm{eff}} > 0),
\]
whereas for $m_{\mathrm{eff}} = 0$ (Bianchi~I):
$\int_0^{\eps} \rho^2\,d\rho/\rho^2 = \eps \to 0$ only after
the additional cancellation that BN23 \cite{BN23} could not achieve
for the massless case.

%% ============================================================
\section{Open Extensions and Companion to BN23}
\label{sec:extensions}
%% ============================================================

\subsection{Bianchi IX (closure of the hierarchy)}

The remaining Type-A Bianchi cosmology, Bianchi~IX ($S^3$ spatial sections),
has $\inf\spec(-\Delta_{S^3}) = 2 > 0$.
The analysis of the present paper applies with minor modifications:
the Plancherel measure is a discrete sum $\sum_{n=1}^\infty n^2$
(spherical harmonics on $S^3$) and the SLE functional becomes
a series rather than an integral.
The spectral gap $= 2$ provides an even stronger IR regulator.

\subsection{Non-compact $H^3$}

For the non-compact $\Gamma=\{e\}$ case, the spectrum of $-\Delta_{H^3}$
is the interval $[1,\infty)$ (purely continuous).
The construction carries through without change;
the compactness assumption on $\Gamma\backslash H^3$ is only needed for the
Plancherel formula to be a genuine $L^2$ isometry.

\subsection{Massive scalar}

For $m>0$, the mode squared-frequency becomes $\omega_\rho^2 = \rho^2 + 1 + m^2 a^2 - 2/\tau^2$,
with even larger spectral gap.
The Hadamard property of SLE for massive fields on Bianchi~V
is a direct corollary of Theorem~\ref{thm:main}.

\subsection{Extensions of the K-L/BN23 program}

\begin{enumerate}
  \item \textbf{Bianchi~V with anisotropy ($a_1 \neq a_2$):}
    The metric \eqref{eq:metric} assumes LRS (locally rotationally symmetric).
    General Bianchi~V has $a_1(t) \neq a_2(t)$ along the $e_1,e_2$ directions.
    The K-L decomposition extends to $H^3$ with anisotropic scale factors
    via the full $SO(3,1)$-representation theory;
    the spectral gap persists since it depends only on the spatial geometry.
  \item \textbf{Beyond dust: stiff matter, radiation.}
    For $w = p/\rho \in [-1,1]$, $a(t) \propto t^{2/(3(1+w))}$.
    The conformal-time potential is $V_{\mathrm{eff}} = (2 - \delta_w)/\tau^2$
    for some $\delta_w \geq 0$; the analysis proceeds identically.
  \item \textbf{Interacting fields.}
    For $\lambda\phi^4$ perturbation theory, the Hadamard state
    constructed here provides the correct initial datum for the
    Brunetti--Fredenhagen--K\"{o}hler renormalization on curved spacetimes
    \cite{BFK95}.
\end{enumerate}

\subsection{The WCH landscape}

In the Wald--Costa--Hollands (WCH) classification of
homogeneous cosmologies by their quantum Hadamard type
(see survey \cite{BN23} \S6),
Bianchi~V matter is the unique remaining \emph{Type-A unconditional} case:
all Bianchi types with $\inf\spec(-\Delta_\Sigma) > 0$ are covered
by the argument of this paper.
The remaining open cases are the non-unimodular types IV, VI, VII with
non-trivial off-diagonal structure.

%% ============================================================
%% References
%% ============================================================

\begin{thebibliography}{99}

\bibitem{BN23}
  R.~Banerjee and M.~Niedermaier,
  \textit{States of Low Energy on Bianchi~I spacetimes},
  J.~Math.~Phys.\ \textbf{64}, 113503 (2023).
  arXiv:2305.11388 [math-ph].
  \textbf{[arXiv-verified.]}

\bibitem{Radzikowski96}
  M.~J.~Radzikowski,
  \textit{Micro-local approach to the Hadamard condition in quantum field theory
  on curved spacetime},
  Commun.~Math.~Phys.\ \textbf{179}, 529--553 (1996).
  DOI: 10.1007/BF02100096.
  \textbf{[Journal-verified via Project Euclid; no arXiv preprint.]}

\bibitem{Wald94}
  R.~M.~Wald,
  \textit{Quantum Field Theory in Curved Spacetime and Black Hole Thermodynamics},
  Chicago University Press (1994).

\bibitem{Olbermann07}
  H.~Olbermann,
  \textit{States of low energy on Robertson--Walker spacetimes},
  Class.~Quantum~Grav.\ \textbf{24}, 5011--5030 (2007).
  arXiv:0704.2986 [gr-qc].

\bibitem{Ellis-MacCallum69}
  G.~F.~R.~Ellis and M.~A.~H.~MacCallum,
  \textit{A class of homogeneous cosmological models},
  Commun.~Math.~Phys.\ \textbf{12}, 108--141 (1969).
  DOI: 10.1007/BF01645908.
  \textbf{[Standard reference for Bianchi classification; no arXiv preprint.]}

\bibitem{Joseph66}
  R.~D.~Joseph,
  \textit{An open universe without a beginning?},
  Phys.~Lett.\ \textbf{20}, 281--282 (1966).
  DOI: 10.1016/0031-9163(66)90537-4.
  (Pre-arXiv journal article; cited for historical attribution only.
  Proposition~\ref{prop:vacuum} provides a self-contained sympy-verified
  proof, so the present paper is logically independent of this
  reference.)

\bibitem{Lebedev65}
  N.~N.~Lebedev,
  \textit{Special Functions and Their Applications},
  Prentice-Hall, Englewood Cliffs, NJ (1965).
  (Revised English translation by R.~A.~Silverman.)
  Internet Archive copy:
  \url{https://archive.org/details/n.-n.-lebedev-special-functions-and-their-applications-1965}.
  (\S5.7: Watson asymptotic for $K_{i\rho}(x)$ at large $\rho$;
  \S5.10: exponential decay $K_\nu(x)\to 0$ as $x\to\infty$;
  \S6.5: K-L Parseval identity \eqref{eq:MS}.)

\bibitem{GK11}
  D.~Goldfeld and A.~Kontorovich,
  \textit{On the determination of the Plancherel measure for
  Lebedev--Whittaker transforms on $GL(n)$},
  (2011).
  arXiv:1102.5086 [math.NT].
  \textbf{[arXiv-verified. Provides elementary K-L inversion proof for $GL(2)$.]}

\bibitem{Helgason84}
  S.~Helgason,
  \textit{Groups and Geometric Analysis},
  Academic Press, New York (1984).
  (Chapters~II--IV: spherical functions and Plancherel for
  symmetric spaces of rank one.)

%% (Faraut 1979 reference removed: replaced by the verified
%%  Helgason \cite{Helgason84} citation throughout the paper.)

\bibitem{DLMF}
  F.~W.~J.~Olver, A.~B.~Olde~Daalhuis, D.~W.~Lozier, B.~I.~Schneider,
  R.~F.~Boisvert, C.~W.~Clark, B.~R.~Miller, B.~V.~Saunders,
  H.~S.~Cohl, and M.~A.~McClain, eds.,
  \textit{NIST Digital Library of Mathematical Functions},
  Release 1.2.4 of 2025-03-15.
  \url{https://dlmf.nist.gov/}.
  (\S10.45: Macdonald function with imaginary order, asymptotic
  expansions for large order.)

\bibitem{Hormander85}
  L.~H\"{o}rmander,
  \textit{The Analysis of Linear Partial Differential Operators IV:
  Fourier Integral Operators},
  Grundlehren der mathematischen Wissenschaften \textbf{275},
  Springer (1985).
  (Theorem~26.1.1: propagation-of-singularities theorem.)

\bibitem{ReedSimonII}
  M.~Reed and B.~Simon,
  \textit{Methods of Modern Mathematical Physics II:
  Fourier Analysis, Self-Adjointness},
  Academic Press (1975).
  (\S X.2: Kato perturbation theory; relatively-bounded perturbations.)

\bibitem{ReedSimonIII}
  M.~Reed and B.~Simon,
  \textit{Methods of Modern Mathematical Physics III: Scattering
  Theory},
  Academic Press (1979).
  (\S XI.3--XI.4: existence and completeness of wave operators
  via Cook's method for $L^{1}$ potentials.)

\bibitem{ParkerToms09}
  L.~Parker and D.~Toms,
  \textit{Quantum Field Theory in Curved Spacetime:
  Quantized Fields and Gravity},
  Cambridge University Press (2009).
  (Chapter~3: mode equations for FLRW cosmologies.)

\bibitem{Ringström09}
  H.~Ringstr\"{o}m,
  \textit{The Cauchy Problem in General Relativity},
  European Mathematical Society (2009).
  (Chapter on Bianchi cosmologies and BKL attractor.)

\bibitem{BKL82}
  V.~A.~Belinskii, I.~M.~Khalatnikov, and E.~M.~Lifshitz,
  \textit{A general solution of the Einstein equations with a time singularity},
  Adv.~Phys.\ \textbf{31}, 639--667 (1982).
  DOI: 10.1080/00018738200101428.
  \textbf{[Standard BKL reference; pre-arXiv.]}

\bibitem{Moretti21}
  V.~Moretti,
  \textit{Fundamental Mathematical Structures of Quantum Theory},
  2nd ed., Springer (2019/2021).
  (Chapter~11: Hadamard states and QFT on curved spacetimes.)

\bibitem{LuedersRoberts90}
  C.~L\"{u}ders and J.~E.~Roberts,
  \textit{Local quasiequivalence and adiabatic vacuum states},
  Commun.~Math.~Phys.\ \textbf{134}, 29--63 (1990).
  DOI: 10.1007/BF02102088.
  \textbf{[Pre-arXiv; journal-accessible. Standard adiabatic vacuum reference.]}

\bibitem{BFK95}
  R.~Brunetti, K.~Fredenhagen, and M.~K\"{o}hler,
  \textit{The microlocal spectrum condition and Wick polynomials
  of free fields on curved spacetimes},
  Commun.~Math.~Phys.\ \textbf{180}, 633--652 (1996).
  arXiv:gr-qc/9510056.
  \textbf{[arXiv-verified.]}

\bibitem{Duistermaat-Guillemin75}
  J.~J.~Duistermaat and V.~W.~Guillemin,
  \textit{The spectrum of positive elliptic operators and periodic bicharacteristics},
  Invent.~Math.\ \textbf{29}, 39--79 (1975).
  \textbf{[Pre-arXiv. Standard WF/propagation-of-singularities reference.]}

\end{thebibliography}

\end{document}
